\section{Results}\label{sec:results}

Figure~\ref{fig:event} traces the licensing difference between upward and downward revisions relative to clean controls.
The three pre-period coefficients are positive but imprecise---0.30, 0.14, and 0.13 log points in years $-4$ through
$-2$---and the joint pre-trend test does not reject equality ($p=0.32$). The difference rises to 0.31 in the revision
year, remains near 0.31 in years $+1$ and $+2$, and reaches 0.57--0.76 in years $+3$ through $+5$. Averaged across event
years 0--5, the licensing gap is 0.495 log points (SE 0.140). In 100,000 direction-label permutations, the two-sided
$p$-value is 0.032.

\begin{figure}[htbp]\centering
\IfFileExists{figures/fig1_event_study_expanded.pdf}{%
\includegraphics[width=0.82\textwidth]{figures/fig1_event_study_expanded.pdf}}{%
\fbox{\parbox{0.78\textwidth}{\centering Run \texttt{python code/build\_expanded\_manuscript\_figures.py} from the repository root to generate this figure.}}}
\caption{Licensing around documentary-reviewed IP-policy revisions}\label{fig:event}
\begin{minipage}{0.9\textwidth}\footnotesize\textit{Notes:} Coefficients are the upward-minus-downward event-time
differences $2\theta_k$ from equation~\eqref{eq:stacked} for the preferred 34-event sample, with institution-clustered
95\% confidence intervals. Year $-1$ is the reference.
\end{minipage}
\end{figure}

\begin{table}[htbp]\centering
\caption{Technology-transfer outcomes around IP-policy revisions}\label{tab:main}
\begin{threeparttable}\footnotesize\setlength{\tabcolsep}{4pt}
\begin{tabular}{lrrrrrrrr}\toprule
 & \multicolumn{5}{c}{Upward minus downward} & \multicolumn{2}{c}{Relative to controls} & \\\cmidrule(lr){2-6}\cmidrule(lr){7-8}
Outcome & Gap & SE & Perm.\ $p$ & BH $q$ & Pre $p$ & Upward & Downward & $N$ \\\midrule
Log licenses and options & 0.495 & (0.140) & 0.032 & 0.159 & 0.32 & 0.279 & -0.216 & 26,555 \\
Filing margin & 0.392 & (0.155) & 0.097 & 0.243 & 0.75 & 0.143 & -0.249 & 26,679 \\
Log new patent applications & 0.260 & (0.175) & 0.205 & 0.341 & 0.89 & 0.136 & -0.124 & 26,679 \\
Log U.S. patents issued & 0.137 & (0.190) & 0.541 & 0.541 & 0.47 & 0.171 & 0.034 & 26,151 \\
Log invention disclosures & -0.132 & (0.095) & 0.380 & 0.475 & 0.40 & -0.007 & 0.124 & 26,679 \\
\bottomrule
\end{tabular}
\begin{tablenotes}[flushleft]\footnotesize
\item \textit{Notes:} Preferred sample: 34 hand-reviewed revisions, eight upward and 26 downward; window $-4$ to $+5$,
year $-1$ omitted. Fixed effects: stack $\times$ institution and stack $\times$ year $\times$ research-size tercile
$\times$ public/private; control: log research expenditure. Permutation $p$: two-sided 100,000-draw reassignment test
preserving eight upward labels. The Monte Carlo SE for the licensing $p$-value is 0.0006. BH $q$: Benjamini--Hochberg
adjustment across the five outcomes. Pre $p$: joint Wald test of the three pre-period direction coefficients.
\end{tablenotes}
\end{threeparttable}
\end{table}

The licensing gap reflects movement in both directions. Relative to non-revising controls, the post-period average is
$+0.279$ after upward revisions and $-0.216$ after downward revisions. The symmetry is not exact, but the result is not
generated solely by one direction. This remains true under the stricter post-support rule: the 29-event sample yields a
gap of 0.499 (SE 0.133), with $+0.230$ for upward and $-0.269$ for downward revisions.

The other outcomes locate rather than generalize the pattern. The filing margin has a positive gap of 0.392 and a flat
pre-period, but its permutation $p$-value is 0.097. New patent applications show a positive gap of 0.260
($p=0.205$). Patents issued have a smaller positive estimate, and invention disclosures show a negative but imprecise
gap. After adjusting the five permutation tests, the licensing result has $q=0.159$. Thus the data do not support a
claim that policy revisions shift the entire commercialization pipeline. The strongest evidence is concentrated in
licensing.

That concentration matters for interpretation. If more supportive communication primarily induced additional faculty
entry into the TTO process, invention disclosures would be a natural first-stage outcome. They are not. Licensing,
however, depends on post-disclosure evaluation, marketing, negotiation, and continued cooperation among inventors, TTO
staff, and firms. The filing margin and applications also lie downstream and move in the same direction, though with
insufficient precision for separate conclusions.

\subsection*{Robustness to the expanded sample definition}

The expansion from 25 to 34 events is not based on the outcome estimates. It follows a complete documentary re-review of
the 37 mechanically eligible pairs that had not previously been hand coded. Nine pass the same C/P and A/B rules.
Nevertheless, late and sparse events could contribute disproportionate leverage. Table~\ref{tab:robust} therefore
shows a uniform stricter outcome-support rule, the reproduced September-25 25-event set, and the broader 59-event
rule-coded specification.

\begin{table}[htbp]\centering
\caption{Alternative event samples: licensing and filing margin}\label{tab:robust}
\footnotesize
\resizebox{0.96\textwidth}{!}{%
\begin{tabular}{lcrrrrrrr}\toprule
 & & \multicolumn{4}{c}{Log licenses and options} & \multicolumn{3}{c}{Filing margin} \\\cmidrule(lr){3-6}\cmidrule(lr){7-9}
Sample & Up/down & Gap & SE & Perm.\ $p$ & Pre $p$ & Gap & SE & Perm.\ $p$ \\\midrule
Preferred documentary sample & 8/26 & 0.495 & (0.140) & 0.032 & 0.32 & 0.392 & (0.155) & 0.097 \\
Strict post-support sample & 7/22 & 0.499 & (0.133) & 0.026 & 0.52 & 0.361 & (0.158) & 0.119 \\
Reproduced September-25 event set & 6/19 & 0.621 & (0.136) & 0.009 & 0.62 & 0.380 & (0.170) & 0.153 \\
Broader sample including rule-coded revisions & 17/42 & 0.419 & (0.113) & 0.006$^{\dagger}$ & 0.34 & 0.187 & (0.104) & 0.176$^{\dagger}$ \\
\bottomrule
\end{tabular}}
\vspace{2pt}

\parbox{0.94\textwidth}{\footnotesize\textit{Notes:} The preferred and strict samples use 100,000 Monte Carlo direction
assignments. The reproduced 25-event row uses exact enumeration over all $\binom{25}{6}=177{,}100$ direction
assignments. The broader 59-event row is retained from the pre-expansion robustness exercise and uses 20,000 draws
($\dagger$); it is not preferred because many pairs are rule-coded rather than document-reviewed.}
\end{table}

The stricter sample is especially informative. It drops all preferred-sample events without at least one observed calendar
year after the event, including both 2023 additions, and also applies the same rule to the original events. The resulting
29-event estimate is nearly identical to the preferred result and has a weaker pre-period difference. The stability is
therefore not an artifact of treating the final AUTM year as a post-revision observation.

No single revision drives the expanded result. Dropping each of the 34 events in turn produces licensing gaps from 0.436
to 0.559 (Appendix Figure~\ref{fig:loo}). The estimate therefore remains positive even when any one of the eight upward
or 26 downward revisions is removed.

\subsection*{Pre-revision balance}

Appendix Table~\ref{tab:balance} compares upward and downward events before revision. Licensing trends are essentially
identical, and differences in pre-revision licensing, research expenditure, public/private status, and timing
classification are not statistically unusual under label permutations. Upward revisions occur somewhat later and at
institutions with more disclosures, but these differences are imprecise. The sharp difference is the starting policy:
mean prior-document PCSI is 0.384 for upward revisions and 0.485 for downward revisions (standardized difference $-1.65$;
permutation $p=0.002$). This mean-reversion concern is substantive and is one reason we do not interpret the direction
contrast causally.

\subsection*{Revision content and supporting annual evidence}

Across all 127 threshold revisions, downward revisions add conflict-of-interest language in 32 percent of cases compared
with 13 percent of upward revisions (Fisher $p=0.019$; Table~\ref{tab:content}). This is a keyword-presence comparison,
not a legal coding of what the provisions require. Provision-level hand coding remains available for the original
25-event content-coded subset (Appendix~\ref{app:provisions}); we retain it as supplementary evidence rather than
retroactively extending those provision labels after observing the expanded-sample results.

\begin{table}[htbp]\centering
\caption{Content added in upward and downward revisions}\label{tab:content}
\begin{threeparttable}\small\begin{tabular}{lrrr}\toprule
 & Upward & Downward & Fisher $p$ \\\midrule
\multicolumn{4}{l}{\textit{Original 25-event content-coded subset} (6 upward, 19 downward)} \\
\quad Adds present assignment language & 33\% & 16\% & 0.562 \\
\quad Adds equity language & 33\% & 47\% & 0.661 \\
\quad Adds conflict of interest language & 0\% & 37\% & 0.137 \\
\quad Adds startup language & 17\% & 11\% & 1.000 \\
\quad Adds copyright language & 33\% & 32\% & 1.000 \\
\multicolumn{4}{l}{\textit{All revisions} (47 upward, 80 downward)} \\
\quad Adds present assignment language & 9\% & 16\% & 0.285 \\
\quad Adds equity language & 28\% & 35\% & 0.437 \\
\quad Adds conflict of interest language & 13\% & 32\% & 0.019 \\
\quad Adds startup language & 15\% & 16\% & 1.000 \\
\quad Adds copyright language & 13\% & 20\% & 0.341 \\
\bottomrule\end{tabular}
\begin{tablenotes}[flushleft]\footnotesize
\item \textit{Notes:} Share of revisions whose later document contains language absent from the earlier document, detected by keyword search: present assignment (``hereby assign''), equity or stock, conflict(s) of interest, startup/spin-off/new venture, copyright. Descriptive; keyword presence does not establish the legal content of a provision.
\end{tablenotes}\end{threeparttable}\end{table}

The annual panel points in the same direction but with much less precision. Table~\ref{tab:twfe} reports regressions of
each outcome on contemporaneous and lagged PCSI. The filing-margin coefficient is positive at every horizon, 0.94 at lag
one and 1.20 at lag two; the licensing coefficient is positive at lags zero to four but imprecise; and disclosure
coefficients are negative at every lag. No lag-one coefficient survives the reported adjustment across the six outcomes.
These annual models are useful as a different representation of the data, but the revision design is preferred because
it imposes documentary timing and comparability requirements directly.

\begin{table}[htbp]\centering
\caption{Supporting evidence: continuous PCSI in two-way fixed-effects models}\label{tab:twfe}
\begin{threeparttable}\footnotesize\setlength{\tabcolsep}{3pt}\begin{tabular}{lrrrrrrr}\toprule
Outcome & Contemp. & Lag 1 & Lag 2 & Lag 3 & Lag 4 & Lag 5 & BH $q$, lag 1 \\\midrule
Invention disclosures & -0.319 & -0.368 & -0.306 & -0.344 & -0.317 & -0.103 & 0.345 \\
 & (0.375) & (0.374) & (0.367) & (0.374) & (0.391) & (0.379) & \\
Filing margin & 0.648 & 0.939 & 1.201 & 0.976 & 0.512 & 0.460 & 0.345 \\
 & (0.527) & (0.505) & (0.520) & (0.468) & (0.445) & (0.497) & \\
New patent applications & 0.329 & 0.571 & 0.895 & 0.632 & 0.196 & 0.357 & 0.345 \\
 & (0.560) & (0.523) & (0.537) & (0.536) & (0.478) & (0.485) & \\
Total patent applications & 0.517 & 0.590 & 0.640 & 0.316 & 0.260 & 0.237 & 0.345 \\
 & (0.574) & (0.525) & (0.499) & (0.488) & (0.499) & (0.469) & \\
Patents issued & 0.267 & 0.548 & 1.205 & 0.981 & 0.774 & 1.300 & 0.345 \\
 & (0.560) & (0.578) & (0.634) & (0.573) & (0.631) & (0.578) & \\
Licenses/options & 0.517 & 0.596 & 0.399 & 0.273 & 0.161 & -0.451 & 0.345 \\
 & (0.489) & (0.452) & (0.405) & (0.441) & (0.467) & (0.503) & \\
\bottomrule\end{tabular}
\begin{tablenotes}[flushleft]\footnotesize
\item \textit{Notes:} Coefficients on PCSI (contemporaneous or lagged $k$ calendar years); institution-clustered standard errors in parentheses. Institution fixed effects on the 149 harmonized institutions, year fixed effects, and lagged log research expenditure, log licensing FTEs and inventor royalty share. Filing-margin samples: $N=2,339, 2,271, 2,132, 2,005, 1,891, 1,772$ from contemporaneous to lag 5. BH $q$: Benjamini--Hochberg adjustment across the six lag-1 outcomes shown.
\end{tablenotes}\end{threeparttable}\end{table}
